\chapter{Une idée pour les formes tordues}

Voici une petite conjecture et ce que je pense en être le début d'une démonstration valable. Le résultat serait une condition suffisante pour que deux modules simples d'évaluation supportés sur une seule orbite aient des extensions non-triviales dans le cas où l'algèbre de courants tordue $\LL$ vérifie la condition $\g^M = \g$ pour tout $M \in \Max S$. C'est le cas pour les formes tordues (voir \cite{LPtfa}, proposition 3.3).

Avant de montrer le résultat, il faut rappeler une notion simple de géométrie algébrique.

\begin{annexeDef} Soit $M \in \Max S$. Alors une \textbf{dérivation de $S$ au point $M$} est une application $k$-linéaire 
$$ \partial : \; S \longrightarrow k $$
qui satisfait à la propriété que pour tout $s_1,s_2 \in S$, on ait 
$$ \partial(s_1s_2) = \partial(s_1) \, s_2(M) + s_1(M) \, \partial(s_2) $$

L'ensemble des dérivations de $S$ au point $M$ est un sous-espace vectoriel de $\Hom_k(S,k)$.
\end{annexeDef}

L'ensemble des dérivations de $S$ au point $M \in \Max S$ est naturellement isomorphe comme espace vectoriel à $(M/M^2)^*$, l'espace tangent associé à $M$. La démonstration est la même que celle donnée dans \cite{GWsym} où on remplace $f \in \CC[t]$ par $s \in S$ et $a \in \CC$ par $M \in \Max S$.

La condition suffisante mentionnée un peu plus haut provient de la proposition qui suit. Elle est inspirée de la proposition 3.4 de l'article \cite{CMspec}.

\begin{annexeConj} \label{trucmargaux} Supposons que l'algèbre de courants tordue $\LL$ satisfasse la propriété que $\g^M = \g$ pour chaque $M \in \Max S$. Soient $V_1$ et $V_2$ des $\g^M$-modules simples de dimension finie. 

Posons $V = V_1^* \otimes V_2$. Alors, il y a une application linéaire injective
\begin{align*}
\beta : \; \Hom_{\g}(\g,V) \otimes (M/M^2)^* \; \longrightarrow \; \cohomo{\LL}{V}
\end{align*}

\em \textbf{Idée de preuve :} Identifions $(M/M^2)^*$ aux dérivations de $S$ au point $M$. Pour ce qui suit, un élément de $(M/M^2)^*$ sera donc une dérivation de $S$ au point $M$. 

L'application $\beta$ envoie un tenseur simple $f \otimes \partial \in \Hom_{\g}(\g,V) \otimes (M/M^2)^*$ sur la restriction à $\LL$ de l'application
\begin{align*}
\tilde{\beta}_{f \otimes \partial} : \quad \g &\otimes S \quad \longrightarrow \quad V\\
x &\otimes s \; \longmapsto \; \partial(s)f(x)
\end{align*}
Notons que cette définition fait de l'application $\beta$ une application $k$-linéaire bien définie. Il faut ensuite montrer que les fonctions linéaires qui composent l'image de $\beta$ sont bien des dérivations de $\LL$ dans $V$. 

Chaque $\tilde{\beta}_{f \otimes \partial}$ est $k$-linéaire par construction. Montrons ensuite que chaque $\tilde{\beta}_{f \otimes \partial}$ est une dérivation de $\g \otimes S$ dans $V$. Soient donc $x \otimes s$ et $y \otimes r \in \g \otimes S$. On peut alors écrire
\begin{align*}
\Big(\tilde{\beta}_{f \otimes \partial}\Big)\big([x \otimes s, y \otimes r]\big) &= \Big(\tilde{\beta}_{f \otimes \partial}\Big)\big([x,y] \otimes sr\big)\\
&= \partial(sr)f\big([x,y]\big)\\
&= \big(\partial(s) \, r(M) + s(M) \, \partial(r)\big)\,f\big([x,y]\big)\\
&= \partial(s) \, r(M)f\big([x,y]\big) + \partial(r) \, s(M)f\big([x,y]\big)\\
&= - \, \partial(s) \, r(M)f\big(y.x\big) + \partial(r) \, s(M)f\big(x.y\big)\\
&= - \, \partial(s) \, r(M)\,y.f\big(x\big) + \partial(r) \, s(M)\,x.f\big(y\big)\\
&= - \, \partial(s) \, (y \otimes r).f(x) + \partial(r) \, (x \otimes s).f(y)\\
&= - (y \otimes r).\big(\partial(s)f(x)\big) + (x \otimes s).\big(\partial(r)f(y)\big)\\
&= (x \otimes s).\big(\partial(r)f(y)\big) - (y \otimes r).\big(\partial(s)f(x)\big)\\
&= (x \otimes s).\Big(\big(\tilde{\beta}_{f \otimes \partial}\big)(y \otimes r)\Big) - (y \otimes r).\Big(\big(\tilde{\beta}_{f \otimes \partial}\big)(x \otimes s)\Big)
\end{align*}
Ainsi, $\tilde{\beta}_{f \otimes \partial} \in \Der(\g \otimes S,V)$ et en conséquence, sa restriction à $\LL \subseteq \g \otimes S$ appartient à l'ensemble de dérivations $\Der(\LL,V)$. En composant avec le passage au quotient $\Der(\LL,V) \rightarrow \cohomo{\LL}{V}$, on obtient bien une application entre les bons ensembles.

Pour terminer, il faudrait montrer que l'application $\beta$ est injective, mais \textbf{je n'ai pas réussi à le faire}. Cependant, je peux montrer que $\tilde{\beta} : \Hom_\g(\g,V) \otimes (M/M^2)^* \rightarrow \cohomo{\g \otimes S}{V}$ est injective. 

Fixons une base $\{\ov{m}_j\}_{j \in J}$ de $(M/M^2)$ et fixons $\{\partial_j\}_{j \in J}$ la base duale associée  à cette dernière. Notons que les $\partial_j$ sont interprétés ici comme des dérivations de l'algèbre $S$ au point $M$. Fixons également une base $\{f_i\}_{i \in I}$ de $\Hom_\g(\g,V)$.

Soit $\sum_{i,j} c_{ij} \, f_i \otimes \partial_j \in \ker \tilde{\beta}$ (la somme étant finie). Alors son image est une dérivation intérieure de $\g \otimes S$ dans $V$. Fixons donc $v \in V$ tel que 
$$ \sum_{i,j} c_{i,j} \, \tilde{\beta}_{f_i \otimes \partial_j} = d^v \in \IDer(\g \otimes S,V)$$ 
Alors, pour tout $x \in \g$, on doit avoir
\begin{align*}
\sum_{i,j} c_{ij} \, \Big(\tilde{\beta}_{f_i \otimes \partial_j}\Big)\big(x \otimes (\textstyle{\sum_k}\, m_k)\big) &= \big(x \otimes (\textstyle{\sum_k}\, m_k)\big).v \\
\sum_{i,j} c_{ij} \, \partial_j\big(\textstyle{\sum_k}\, m_k\big) f_i(x) &= \big(\textstyle{\sum_k}\, m_k\big)(M) \; x.v \\
\sum_{i,j} c_{ij} \, f_i(x) &= 0 \\
\sum_i \Big(\sum_j c_{ij}\Big) \, f_i(x) &= 0
\end{align*}
Puisque les $\{f_i\}_{i \in I}$ forment une base de $\Hom_\g(\g,V)$, on obtient que $\sum_j c_{i,j} = 0$ pour chacun des indices $i$. Alors on conclut que
$$ \sum_{i,j} c_{ij} \, f_i \otimes \partial_j = \sum_i \Big(\sum_j c_{ij}\Big) \, f_i \otimes \partial_j = 0 \in \Hom_\g(\g,V) \otimes (M/M^2)^*$$
Ainsi, $\tilde{\beta}$ est injective. 

Le problème avec cette méthode pour prouver l'injectivité de $\beta$ est que les éléments de la forme $x \otimes (\sum_k m_k)$, où $x \in \g$ est arbitraire, ne sont pas forcément dans $\LL$ alors on ne peut pas utiliser ces éléments pour aboutir à la conclusion analogue à celle ci-haut dans le cas de $\g \otimes S$. \flushright \textsc{Fin de l'idée de preuve.}
\end{annexeConj}

\begin{annexeRem} La preuve de l'injectivité de $\tilde{\beta}$ ci-haut généralise la construction de V.Chari et A.Moura donnée à la proposition 3.4 de leur article \cite{CMspec}. La différence est qu'ici, l'algèbre de Lie en question est $\g \otimes S$ au lieu de $\g \otimes k[t^{\pm 1}]$.
\end{annexeRem}


\begin{annexeConj} Soit $M \in \Max S$ et soient $V_1$ et $V_2$ des $\g^M$-modules simples de dimension finie. Posons $V = V_1^* \otimes V_2$. Alors
$$ \left.\begin{array}{c} \Hom_{\g^M}(\g^M,V) \neq \{0\} \\ (M/M^2)^* \neq \{0\}\end{array}\right] \quad \Longrightarrow \quad \Ext^1_\LL(V_1,V_2) \neq \{0\} $$
\end{annexeConj}

\begin{annexeRem} Ces deux conjectures donneraient un sens à la possible « interprétation géométrique » à laquelle E.Neher et A.Savage faisaient référence suite à leur lemme 4.5 dans \cite{NSext}. Voir la remarque 4.6 de ce même article.
\end{annexeRem}